
We propose DPitG as a sequential stopping rule that jointly requires posterior
precision and a conclusive verdict before data collection ends.
PitG eliminates the early-peeking bias of coupled methods by decoupling the stopping
and decision criteria \citep{kruschke2015doing}, but accepts high inconclusive rates
as the price. Inconclusive rates rise steeply the closer the null hypothesis is to the true value.

DPitG closes this gap by imposing both requirements simultaneously: data collection
continues until $\omega_{\rm HDI} \le \omega_{\rm goal}$ \textit{and} the posterior's
position relative to the ROPE yields a conclusive verdict.
This joint requirement delivers a $2.6\times$ gain in conclusiveness over PitG at a
median cost of only 4.7\% (4\% among conclusive experiments) more samples
($\theta_{\rm true}=0.5$, $\omega_{\rm goal}=0.08$,
${\rm ROPE}_{\rm min}^{\rm max}=[0.45,0.55]$),
while the HDI+ROPE algorithm accrues a systematic false-positive rate because its
coupled design can stop on wide, early posteriors before they are representative.
The gain is robust: DPitG's conclusiveness holds nearly flat ($97.3\%$--$98.0\%$)
across the same $\omega_{\rm goal}$ range, over which PitG's conclusiveness
varies widely\footnote{And collapses to zero when $\omega_{\rm goal} = {\rm ROPE}_{\rm max} - {\rm ROPE}_{\rm min}$.}
and is largest near the ROPE boundaries where evidence is hardest to read.

Unlike NHST, which can only reject or fail to reject the null (see Appendix~\ref{app:nhst}),
the HDI+ROPE framework's decision rule enables positive \textit{acceptance} of the null
when the entire HDI falls within the ROPE, a critical advantage for confirmatory
research areas such as bioequivalence, safety testing, and software regression testing;
DPitG reliably delivers this verdict.

Prospective planning is straightforward: the closed-form
$N_{\rm goal}(\theta_{\rm null}, \omega_{\rm goal}) \propto
\theta_{\rm null}(1-\theta_{\rm null})/\omega_{\rm goal}^2$
provides the minimum budget required, with $N_{\rm max}$ set generously above this
baseline to ensure DPitG reaches its full conclusiveness potential.

The framework extends to single-group continuous data
(Appendix~\ref{app:expected_stop}) and to two-group comparisons, both binomial
and continuous (Appendix~\ref{app:extensions}).\footnote{These extensions are
theoretical derivations; stopping-property validation beyond the single-group Bernoulli
case remains to be completed in a dedicated simulation study.}
All algorithms are implemented in openly available
code\footnote{\href{https://github.com/elzurdo/dpitg}{https://github.com/elzurdo/dpitg}},
with an interactive online calculator for prospective sample-size planning and
retrospective analysis available at
\href{https://r-w-t-y.streamlit.app}{https://r-w-t-y.streamlit.app}.

DPitG is the method of choice whenever a reliable verdict (acceptance, rejection,
or inconclusive) is required and the budget is sufficient to reach
$N_{\rm goal}(\theta_{\rm null})$; the closed-form $N_{\rm goal}$ makes this judgement
transparent before any data are collected.
